\chapter{The ingestion pipeline}
\label{ch:ingestion}

Ingestion is the part of the system everything else depends on, so it is
deliberately the least clever part. One call to \code{IngestService.ingest()}
takes a file and leaves behind a document row, its pages, its chunks and their
embeddings --- or raises, having stored nothing but the file copy.

\begin{shellblock}
file --> hash --> duplicate?  --yes--> return {"status": "duplicate"}
                     |no
                     v
                  store copy in data/documents/
                     |
                     v
                  parse  (pdf | docx | pptx)  -> list[ParsedSection]
                     |
                     v
                  selective OCR (PDF only, sparse pages)
                     |
                     v
                  normalize -> list[CleanSection]   (drops empty pages)
                     |
                     v
                  chunk     -> list[ChunkPayload]   (650 words / 75 overlap)
                     |
                     v
                  persist: Document + DocumentPage* + DocumentChunk*
                     |
                     v
                  embed each chunk -> ChunkEmbedding (float32 blob)
                     |
                     v
                  commit + invalidate the vector index
\end{shellblock}

\section{Supported formats}

\file{.pdf}, \file{.docx} and \file{.pptx}. Anything else raises
\code{ValueError} from the service and becomes \code{400 Unsupported file type}
at the HTTP layer.

\section{Deduplication}

The file is hashed (SHA-256, streamed) before anything else. If a document with
that hash exists, ingestion stops and returns \code{status: "duplicate"} with
the existing document id --- no second copy, no re-parse. \code{file\_hash}
carries a unique index, so the invariant is enforced by the database and not
only by the check.

This is exact-bytes deduplication. Two exports of the same report with different
metadata are different files here; near-duplicates are a separate concern
handled in Section~\ref{sec:duplicate-detection}.

\section{Storage}

The accepted file is copied into \code{<DATA\_DIR>/documents} under
\code{<sanitized stem>\_\_<first 8 hex of the hash><extension>} --- readable
enough to recognize in a file listing, unique enough not to collide --- and
\code{documents.file\_path} records where it went.
\code{resolve\_document\_file\_path()} re-resolves that path at read time, so an
installation that was moved to another drive still finds its files: it tries the
stored path, then the path relative to the current documents directory, then the
bare file name inside it.

\section{Parsers}

Every parser returns the same structure --- \code{list[ParsedSection]} with a
page number, the raw text, an optional section title, and the extraction
provenance fields. Weak documents produce empty-safe output rather than
exceptions.

\begin{description}
\item[PDF (\file{app/parsers/pdf\_parser.py})] \code{pypdf} extracts text page by
page. A section title is guessed from the first line that is either all-caps or
ends with a colon. After extraction the page list is handed to the selective OCR
service, and any page that came back from OCR gets its title re-detected from
the new text.

\item[DOCX (\file{app/parsers/docx\_parser.py})] \code{python-docx} walks
paragraphs. A paragraph whose style name starts with \code{Heading} closes the
current section and opens a new one, so ``pages'' here are logical sections
numbered from 1. Table rows are appended as pipe-separated lines, which keeps
tabular values searchable without pretending to preserve layout.

\item[PPTX (\file{app/parsers/pptx\_parser.py})] Slides are read straight from
the OOXML inside the \file{.pptx} archive with the standard library --- no
PowerPoint and no extra dependency. Slide order comes from the numeric part of
\file{ppt/slides/slideN.xml}; each slide is one section, and its first
meaningful line becomes the section title.
\end{description}

\section{Selective OCR}

\code{SelectiveOCRService} exists to rescue scanned pages without paying for OCR
on documents that do not need it. It applies only to PDFs, and only to pages
whose native text scores below \env{OCR\_MIN\_TEXT\_CHARACTERS}.

\begin{itemize}
\item The page is rasterized with \code{pymupdf} at \env{OCR\_DPI} and passed to
Tesseract with \env{OCR\_LANGUAGES}.
\item Native and OCR text are merged rather than swapped, so a page that had a
little real text does not lose it.
\item The section records \code{extraction\_method="ocr"} and
\code{ocr\_attempted=True}, which is what makes the result auditable afterwards.
\item If Tesseract, its language data or \code{pymupdf} is missing, the service
reports itself unavailable and ingestion continues with native text only.
\end{itemize}

\section{Normalization}

\code{normalize\_sections()} (\file{app/processing/text\_cleaner.py}) does three
things.

\begin{enumerate}
\item \textbf{Clean each page.} NFC normalization, \code{\textbackslash r\textbackslash n}
folding, per-line whitespace collapse, and at most one blank line between
blocks. Turkish characters are preserved exactly --- folding for \emph{matching}
happens later and separately, in \file{app/text/normalize.py}.

\item \textbf{Drop repeated page artefacts.} A line 3--120 characters long that
appears on at least half the pages (minimum two) is treated as a running header
or footer and removed from every page.

\item \textbf{Drop pages that cleaned to nothing.} Their page numbers survive
only in the extraction quality rollup, as \code{empty\_pages}.
\end{enumerate}

The cleaning is intentionally conservative. Aggressive cleaning removes wording
that retrieval needs; both the raw and the cleaned text are stored, so nothing
is lost either way.

\section{Chunking}

\code{chunk\_sections()} walks the cleaned sections and emits overlapping word
blocks.

\begin{longtable}{L{0.24\textwidth} L{0.14\textwidth} L{0.50\textwidth}}
\toprule
\textbf{Parameter} & \textbf{Default} & \textbf{Note} \\
\midrule
\endfirsthead
\bottomrule
\endfoot
\code{target\_words} & \code{650} & Within the 500--800 range the project settled
on. \\
\code{overlap\_words} & \code{75} & Must be smaller than \code{target\_words};
otherwise \code{ValueError}. \\
\end{longtable}

Each chunk carries \code{chunk\_order} (global, 1-based, in document order), the
page range it came from, and the section title of its source section. Chunks are
the searchable unit: keyword search, semantic search, review rules and
comparison all read chunks, never raw pages.

\section{Extraction quality}

\file{app/processing/extraction\_metrics.py} computes per-page metrics and rolls
them up into \code{documents.extraction\_quality} (a JSON column). This is the
field to read when an answer looks thin --- it usually means the document
arrived as images.

\begin{longtable}{L{0.32\textwidth} L{0.58\textwidth}}
\toprule
\textbf{Key} & \textbf{Meaning} \\
\midrule
\endfirsthead
\bottomrule
\endfoot
\code{page\_count} & Pages the parser produced. \\
\code{stored\_page\_count} & Pages that survived normalization and have a
\code{document\_pages} row. \\
\code{ocr\_page\_count} / \code{ocr\_attempted\_page\_count} & Pages whose text
came from OCR, and pages OCR was tried on. \\
\code{sparse\_page\_count} & Pages with text below the sparse threshold. \\
\code{no\_text\_page\_count} / \code{no\_text\_pages} & Unreadable pages,
including the ones normalization dropped. \\
\code{empty\_pages} & Page numbers dropped at normalization. They have no page
row, so this is the only record they exist at all. \\
\code{total\_chars} / \code{mean\_chars\_per\_page} & Size of what was
extracted. \\
\end{longtable}

The value is \code{NULL} for documents ingested before the column existed; treat
that as \emph{unknown}, not as \emph{perfect}.

\section{Embeddings}

Each chunk is embedded with the active provider and stored as packed
little-endian \code{float32} bytes in \code{chunk\_embeddings}. A vector with no
signal --- the zero vector a \code{token-hash} provider returns for text with no
usable tokens --- is skipped rather than stored, so \code{embeddings\_created}
can be lower than the chunk count.

After the commit, \code{invalidate\_vector\_index()} drops the cached matrix so
the next search sees the new document.

\section{Ingestion endpoints}

\begin{description}
\item[\route{POST /ingest}] One file, one synchronous response. Returns document
id, status, page and chunk counts, \code{ocr\_pages}, the extraction quality
rollup, the number of embeddings created and the provider name.

\item[\route{POST /ingest/batch}] Many files. Uploads are staged to temporary
files first --- an \code{UploadFile} is gone once the request ends --- and the
work is handed to a background job, which returns \code{202} with a
\code{job\_id}. Each file gets its own database session, so one failure does not
poison the rest of the batch; failures land in the per-item result with
\code{status: "error"} and a message.
\end{description}

The catalog ingestion endpoints (Chapter~\ref{ch:catalog}) reuse the same
service to pull files off network shares.

\section{Result payload}

\begin{jsonblock}
{
  "document_id": 42,
  "status": "ingested",
  "file_name": "2025-BIG-E-NVH-01.pdf",
  "pages": 18,
  "ocr_pages": 2,
  "extraction_quality": {"...": "..."},
  "chunks": 57,
  "embeddings_created": 57,
  "embedding_provider": "token-hash-v1"
}
\end{jsonblock}
